\subsection{Introduction to the FJ model}
\label{subsec:fj-intro}

\subsubsection{The FJ model}
\label{subsubsec:fj-model}

In the FDG model every agent keeps averaging the opinions it listens to, and a
persistent disagreement can only come from the structure of the network, i.e.\ from
several closed groups (Theorem~\ref{thm:disagreement}). The Friedkin--Johnsen (FJ)
model adds a second ingredient: every agent is also tied to a fixed opinion of its own,
called its \emph{prejudice}. Following~\cite{proskurnikov2017tutorial}, the model reads
\begin{equation}\label{eq:fj}
  x(k+1) = \Lambda W x(k) + (I - \Lambda)\, u ,
\end{equation}
where
\begin{itemize}
  \item $W \in \mathbb{R}^{n \times n}$ is a row-stochastic matrix, the \emph{matrix of
        social influences}; with the convention of Section~\ref{subsec:prelim},
        $w_{ij} > 0$ means that agent $i$ is influenced by agent $j$, so $W$ plays the
        role that $A$ plays in the FDG model;
  \item $\Lambda = \operatorname{diag}(\lambda_1, \dots, \lambda_n)$, with
        $\lambda_i \in [0,1]$ the \emph{susceptibility} of agent $i$ to social
        influence;
  \item $u \in \mathbb{R}^n$ is a constant vector collecting the agents'
        \emph{prejudices}.
\end{itemize}
Componentwise,
\begin{equation}\label{eq:fj-i}
  x_i(k+1) = \lambda_i \sum_{j=1}^{n} w_{ij}\, x_j(k) + (1 - \lambda_i)\, u_i .
\end{equation}
At every step agent $i$ averages the opinions it listens to, exactly as in the FDG
model, and mixes this average with its prejudice: the weight $\lambda_i$ goes to the
others and the complement $1 - \lambda_i$ goes to $u_i$. The two quantities play
different roles: the prejudice $u_i$ is an opinion and says \emph{where} the agent is
anchored, while $1 - \lambda_i$ is a weight and says \emph{how strongly}. The same
model appears in~\cite{bullo2024lns}, where $\lambda_i$ is called the openness level of
agent $i$ and $1 - \lambda_i$ its attachment to the prejudice.

The special cases of~\eqref{eq:fj} are the following~\cite{proskurnikov2017tutorial}.
\begin{itemize}
  \item If $\lambda_i = 1$, agent $i$ follows the FDG rule. If $\Lambda = I$, the FJ
        model reduces to the FDG model with $A = W$; for this reason we assume
        $\Lambda \neq I$ from now on.
  \item If $\lambda_i < 1$, agent $i$ is called \emph{prejudiced}: it is anchored to its
        prejudice $u_i$.
  \item If $\lambda_i = 0$, then $x_i(k) = u_i$ for all $k \geq 1$: the agent ignores
        the others and jumps to its prejudice after one step. Its opinion never changes,
        i.e.\ the agent is stubborn in the sense of Definition~\ref{def:stubborn}, only
        if in addition $u_i = x_i(0)$. A stubborn agent of the FJ model is therefore an
        agent with $\lambda_i = 0$ and $u_i = x_i(0)$.
\end{itemize}

In social terms, the self-weight of the FDG model and the prejudice of the FJ model are
two different forms of resistance to influence. A self-weight ties an agent to its
\emph{current} opinion, which still moves towards the others
(Section~\ref{subsubsec:fdg-stub-model}); a prejudice ties it to a \emph{fixed}
opinion, which never moves.

In the original formulation of Friedkin and Johnsen the prejudices coincide with the
initial opinions, $u = x(0)$, and the susceptibilities are coupled to the self-weights
through $1 - \lambda_i = w_{ii}$~\cite{proskurnikov2017tutorial,bullo2024lns}. As
in~\cite{proskurnikov2017tutorial}, we do not impose these conditions: $u$ and $x(0)$
are independent, and so are $\Lambda$ and $W$. The choice $u = x(0)$ will be needed in
Section~\ref{subsec:fj-stub} to define the influence centrality.

\subsubsection{Convergence and stability}
\label{subsubsec:fj-conv}
% RESTRUCTURED (2026-10-23): this subsubsection was called "Convergence: all agents
% P-dependent" and started from Corollary 22 of the tutorial. It now follows the order
% of the tutorial: definitions, P-dependent agents and decomposition, Theorem 21 (where
% V is defined), Corollary 22 (all agents P-dependent), Corollary 23, and only at the
% end our observations (error dynamics). The former subsubsection "The general case:
% P-independent agents" (label subsubsec:fj-general) is merged here.

The FJ model~\eqref{eq:fj} is a linear system with the constant input $u$. Following
\cite{proskurnikov2017tutorial}, we say that it \emph{converges} if the limit
$x(\infty) = \lim_{k \to +\infty} x(k)$ exists for every initial condition $x(0)$, and
that it is \emph{asymptotically stable} if $\Lambda W$ is convergent, i.e.\
$\rho(\Lambda W) < 1$ (Definition~\ref{def:semiconv} and
Theorem~\ref{thm:spectral-radius}). Following the same reference, a matrix $A$ is
called \emph{regular} if $\lim_{k \to +\infty} A^k$ exists, i.e.\ if it is
semi-convergent in the sense of Definition~\ref{def:semiconv}.

Both properties are characterized by distinguishing the agents that are reached by the
prejudices from those that are not.

\begin{definition}[P-dependent agents~\cite{proskurnikov2017tutorial}]
\label{def:p-dependent}
An agent is \emph{P-dependent} if it is prejudiced, or if it is influenced, directly or
through intermediaries, by some prejudiced agent, i.e.\ if there is a directed walk in
$G$ from the agent to a prejudiced agent. Otherwise, the agent is
\emph{P-independent}.
\end{definition}

The walk goes from the agent to the prejudiced agent because of our convention on the
edges; in~\cite{proskurnikov2017tutorial}, where the edges are drawn in the opposite
direction, the same walk goes from the prejudiced agent to the agent. In social terms,
an agent is P-dependent if at least one prejudice reaches it.

A P-independent agent $i$ is not prejudiced, so $\lambda_i = 1$, and it is not
influenced by any P-dependent agent: if $w_{ij} > 0$ with $j$ P-dependent, there would
be a walk from $i$ to a prejudiced agent through $j$, and $i$ would be P-dependent.
Since $\Lambda \neq I$, there is at least one prejudiced agent, hence at least one
P-dependent agent. Renumbering the agents so that agents $1, \dots, r$, with
$r \geq 1$, are P-dependent and agents $r+1, \dots, n$ are P-independent, and
splitting $x = [x^1; x^2]$, $u = [u^1; u^2]$, $W$ and $\Lambda$ accordingly, we
have $W^{21} = 0$ and $\Lambda^{22} = I$, and the model~\eqref{eq:fj} decomposes
as~\cite{proskurnikov2017tutorial}
\begin{align}
  x^1(k+1) &= \Lambda^{11}\bigl[W^{11} x^1(k) + W^{12} x^2(k)\bigr]
              + (I - \Lambda^{11})\, u^1, \label{eq:fj-dep}\\
  x^2(k+1) &= W^{22} x^2(k), \label{eq:fj-indep}
\end{align}
where $W^{22}$ is row-stochastic, since $W$ is row-stochastic and $W^{21} = 0$. The
P-independent agents therefore follow an FDG
model among themselves, with $A = W^{22}$, and no prejudice reaches them; the
P-dependent agents receive the opinions of the P-independent ones as an additional
input.

\begin{theorem}[\cite{proskurnikov2017tutorial}]
\label{thm:fj-general}
Let the group have $r \geq 1$ P-dependent agents and $n - r \geq 0$ P-independent
ones. Then:
\begin{enumerate}
  \item[(i)] the subsystem~\eqref{eq:fj-dep} is asymptotically stable, i.e.\
        $\rho(\Lambda^{11} W^{11}) < 1$;
  \item[(ii)] the model~\eqref{eq:fj} converges if and only if $r = n$ or $W^{22}$ is
        regular;
  \item[(iii)] in this case
        \begin{equation}\label{eq:fj-V}
          x^1(\infty) = V \begin{bmatrix} u^1 \\ x^2(\infty) \end{bmatrix}, \qquad
          V = (I - \Lambda^{11} W^{11})^{-1}
              \begin{bmatrix} I - \Lambda^{11} & \Lambda^{11} W^{12} \end{bmatrix},
        \end{equation}
        and the matrix $V$ is row-stochastic.
\end{enumerate}
\end{theorem}

Since $V$ is row-stochastic, the final opinion of every P-dependent agent is a convex
combination of the prejudices and of the final opinions of the P-independent
agents~\cite{proskurnikov2017tutorial}. The latter are the limit of the FDG
model~\eqref{eq:fj-indep}: by the Rosetta stone (Section~\ref{subsubsec:rosetta}),
$W^{22}$ is regular if and only if every sink of the condensation digraph of the
subgraph of P-independent agents, whose adjacency matrix is $W^{22}$, is aperiodic. In
social terms, the P-dependent agents settle on a compromise between the prejudices and
the opinions of the P-independent agents they listen to whenever these opinions
settle; whether the whole group settles is decided only by the P-independent agents,
who ignore the prejudices and behave as in the FDG model.

When $r = n$ there are no P-independent agents, the subsystem~\eqref{eq:fj-dep} is the
whole model, and~\eqref{eq:fj-V} reduces to
\begin{equation}\label{eq:fj-V-all}
  x(\infty) = V u, \qquad V = (I - \Lambda W)^{-1}(I - \Lambda),
\end{equation}
now with $V \in \mathbb{R}^{n \times n}$. This case is the subject of the following
result.

\begin{corollary}[\cite{proskurnikov2017tutorial}]
\label{cor:fj-stable}
The FJ model~\eqref{eq:fj} is asymptotically stable if and only if all agents are
P-dependent. Then $V = (I - \Lambda W)^{-1}(I - \Lambda)$ and $x(\infty) = Vu$ for
every $x(0)$. This holds, in particular, if $\lambda_i < 1$ for every agent, or if
$\Lambda \neq I$ and $G$ is strongly connected ($W$ irreducible).
\end{corollary}

In social terms, the group converges, for every initial opinion, as soon as every agent
that is not anchored to a prejudice listens, directly or through intermediaries, to
someone who is. On a strongly connected network, such as the two networks of
Figure~\ref{fig:networks}, a single prejudiced agent is enough. Note that this does not
require aperiodicity, unlike the FDG model of Section~\ref{subsec:fdg-stub}.

The condition of Corollary~\ref{cor:fj-stable} can also be read on the matrix
$\Lambda W$ with the tools of the course. The matrix $\Lambda W$ is no longer
row-stochastic: since $W$ is row-stochastic, the $i$-th row of $\Lambda W$ sums to
$\lambda_i$. Rows of agents that are not prejudiced sum to $1$, rows of prejudiced
agents sum to less than $1$, and since $\Lambda \neq I$ at least one such row exists.
Hence $\Lambda W$ is row-substochastic (Definition~\ref{def:row-substochastic}) and
Lemma~\ref{lem:substochastic} applies.
% MOVED (2026-10-23): Definition "Row-substochastic matrix" and Lemma "Convergent
% row-substochastic matrices" were here; they are now in the Preliminaries
% (Section 1.1.2), with the definition of weighted out-degree.

For $A = \Lambda W$ the weighted out-degree of agent $i$ is $\lambda_i$: the nodes with
out-degree $1$ are the agents that are not prejudiced, and the nodes with out-degree
less than $1$ are the prejudiced agents. The condition of
Lemma~\ref{lem:substochastic} therefore reads: every agent that is not prejudiced is
influenced, directly or through intermediaries, by a prejudiced agent, i.e.\ all the
agents are P-dependent. This is exactly the condition of
Corollary~\ref{cor:fj-stable}. Indeed,~\cite{proskurnikov2017tutorial} remarks that
Corollary~\ref{cor:fj-stable} can be turned into a criterion of Schur stability for
substochastic matrices, and~\cite{bullo2024lns} proves the same condition for the FJ
model through this lemma.

If some agents are P-independent, the model is not asymptotically stable, but it may
still converge. Part (ii) of Theorem~\ref{thm:fj-general} gives the condition in terms
of $W^{22}$; in terms of the whole matrix $\Lambda W$ it reads as follows.

\begin{corollary}[\cite{proskurnikov2017tutorial}]
\label{cor:fj-regular}
The FJ model~\eqref{eq:fj} converges if and only if $\Lambda W$ is regular.
\end{corollary}

As remarked in~\cite{proskurnikov2017tutorial}, this property is not trivial: a system
$x(k+1) = A x(k) + B u$ with a regular matrix $A$ may not converge, and may even have
unbounded solutions, as for $A = B = I$.

We conclude with two observations that are not stated
in~\cite{proskurnikov2017tutorial}: the first one relies on~\cite{bullo2024lns}, the
second one is an intuition of ours.

First, the limit of Corollary~\ref{cor:fj-stable} is the unique equilibrium of the
model. If $\Lambda W$ is convergent, \eqref{eq:fj} is an affine system
$x(k+1) = A x(k) + b$ with $A = \Lambda W$ convergent and $b = (I - \Lambda)u$; then
$I - \Lambda W$ is invertible, the only equilibrium is
$x^* = (I - \Lambda W)^{-1}(I - \Lambda)u = Vu$, and $x(k) \to x^*$ for every
$x(0)$~\cite{bullo2024lns}. Letting $e(k) = x(k) - x^*$ and subtracting
$x^* = \Lambda W x^* + (I - \Lambda)u$ from~\eqref{eq:fj} gives
\begin{equation}\label{eq:fj-error}
  e(k+1) = \Lambda W\, e(k),
\end{equation}
an autonomous linear system in which the prejudice $u$ does not appear: $\Lambda W$
decides whether the opinions converge, $u$ decides where.

Second, the linear system~\eqref{eq:fj-error} gives an intuition for the role of the
P-dependent agents. At each step, agent $i$ keeps a fraction $\lambda_i$ of the
weighted average of the deviations of the agents it listens to. An agent that is not
prejudiced ($\lambda_i = 1$) only averages the deviations, while a prejudiced agent
loses a fraction $1-\lambda_i$ of them (in~\eqref{eq:fj}, this fraction is replaced by
the agent's fixed prejudice). The deviations therefore leave the group only through
the prejudiced agents, and they vanish only if every agent is influenced, directly or
through intermediaries, by one of them, i.e.\ all agents are P-dependent. A group of
agents that is not influenced, even indirectly, by any prejudiced agent loses no
deviation and behaves as an FDG model, as in~\eqref{eq:fj-indep}.

\subsubsection{The equilibrium and the augmented graph}
\label{subsubsec:fj-equilibrium}

Assume that all agents are P-dependent, so that the opinions converge to
$x(\infty) = Vu$, with $V = (I - \Lambda W)^{-1}(I - \Lambda)$
(Corollary~\ref{cor:fj-stable}). By Theorem~\ref{thm:fj-general}, the matrix $V$ is
row-stochastic; see also~\cite{bullo2024lns}. The row sums can also be checked
directly: since $W \mathbf{1}_n = \mathbf{1}_n$, we have
$(I - \Lambda W)\mathbf{1}_n = (I - \Lambda)\mathbf{1}_n$, and multiplying on the left
by $(I - \Lambda W)^{-1}$ gives $V \mathbf{1}_n = \mathbf{1}_n$.

Hence each final opinion
\[
  x_i(\infty) = \sum_{j=1}^{n} v_{ij}\, u_j
\]
is a weighted average of the prejudices, and three facts follow.
\begin{itemize}
  \item The initial opinions are forgotten: $x(0)$ does not appear in $x(\infty)$.
  \item Only the prejudices of prejudiced agents count: if $\lambda_j = 1$, the $j$-th
        column of $I - \Lambda$ is zero, hence so is the $j$-th column of $V$, and
        $u_j$ has no weight in any final opinion.
  \item Every final opinion lies between the smallest and the largest prejudice of the
        prejudiced agents.
\end{itemize}
In the FDG model, the consensus value $w^\top x(0)$ is a weighted average of the
initial opinions, with the same weights $w$ for every agent
(Theorem~\ref{thm:consensus}). In the FJ model, the final opinions are weighted
averages of the prejudices, and each agent has its own weights, given by its row of
$V$.

% REMOVED (2026-10-23): \subsubsection{The general case: P-independent agents}
% \label{subsubsec:fj-general} -- merged into Section 3.1.2 (Theorem 21 and
% Corollary 23 of the tutorial), to follow the order of the tutorial.
